\documentclass[11pt,a4paper]{article}
\usepackage[utf8]{inputenc}
\usepackage[T1]{fontenc}
\usepackage{amsmath,amsfonts,amssymb}
\usepackage{graphicx}
\usepackage{hyperref}
\usepackage{listings}
\usepackage{color}
\usepackage{booktabs}
\usepackage{float}
\usepackage{algorithm}
\usepackage{algorithmic}
\usepackage{geometry}
\geometry{margin=1in}
\definecolor{zenblue}{RGB}{41,121,255}
\hypersetup{colorlinks=true,linkcolor=zenblue,urlcolor=zenblue,citecolor=zenblue}

\title{\textbf{Mathematical Reasoning in Zen Models}\\
\large Technical Report v2025.08}
\author{Zen LM Research Team\\
\texttt{research@zenlm.org}}
\date{August 2025}

\begin{document}
\maketitle

\begin{abstract}
We present a comprehensive study of mathematical reasoning capabilities in Zen models,
covering training data construction, process reward models (PRMs) for step-level
supervision, and a symbolic-neural hybrid verification system. Zen achieves 73.2\%
on the MATH benchmark, 96.4\% on GSM8K, 42.1\% on AIME 2024, 78.4\% on AMC 2023,
and 31.2\% on FrontierMath. Our step-level PRM improves MATH accuracy by 8.1 points
over outcome-only reward models. The hybrid verification system using SymPy for
algebraic verification and neural methods for proof checking enables reliable
answer verification across problem classes. We describe data construction, training
protocol, and detailed performance analysis by mathematical domain.
\end{abstract}

\section{Introduction}

Mathematical reasoning requires multi-step inference, precise computation, and
systematic verification -- capabilities that push language models to their limits.
Unlike open-ended generation tasks, mathematics has ground truth: answers are
either correct or incorrect, enabling automated evaluation and scalable training
signal generation.

The Zen mathematical reasoning stack consists of three components:
\begin{enumerate}
  \item \textbf{Mathematical pretraining data}: 120B tokens of mathematical text
    spanning textbooks, papers, and competition problems.
  \item \textbf{Process Reward Model (PRM)}: A model that evaluates solution quality
    step-by-step rather than only at the final answer.
  \item \textbf{Hybrid verification}: Symbolic algebra (SymPy) + neural verification
    for reliable answer checking.
\end{enumerate}

\section{Mathematical Pretraining Data}

\subsection{Corpus Composition}

\begin{table}[H]
\centering
\begin{tabular}{lrr}
\toprule
\textbf{Source} & \textbf{Tokens} & \textbf{Quality} \\
\midrule
arXiv math papers & 38B & High \\
Math textbooks (OCR) & 22B & High \\
Competition problems + solutions & 4B & Very high \\
Stack Exchange (math) & 18B & Medium \\
Synthetic generated problems & 24B & High \\
Wikipedia math articles & 8B & High \\
Code (mathematical libraries) & 6B & Medium \\
\midrule
Total & 120B & -- \\
\bottomrule
\end{tabular}
\caption{Mathematical pretraining corpus.}
\end{table}

\subsection{Synthetic Problem Generation}

A core challenge is data scarcity at the frontier: competition-level problems are
rare and quickly memorized. We generate synthetic problems via three strategies:

\subsubsection{Template Perturbation}

Existing problems are modified by changing numerical constants, function types, or
geometric configurations while preserving problem structure. For each seed problem,
we generate 10--20 variants with verified solutions.

\subsubsection{Forward-Generation}

We sample mathematical objects (functions, sequences, geometric configurations) and
generate problems whose answer is the sampled object. For example: sample a polynomial
$p(x)$, generate a word problem whose solution is $p(x)$.

\subsubsection{Backward-Chaining}

Starting from a target theorem or formula, generate a minimal problem that requires
the target as a solution step. This biases the distribution toward problems that
exercise specific mathematical tools.

\begin{table}[H]
\centering
\begin{tabular}{lrr}
\toprule
\textbf{Strategy} & \textbf{Problems Generated} & \textbf{Verified Correct} \\
\midrule
Template perturbation & 4.2M & 97.3\% \\
Forward-generation & 2.8M & 94.1\% \\
Backward-chaining & 1.6M & 91.8\% \\
Total & 8.6M & 94.9\% \\
\bottomrule
\end{tabular}
\caption{Synthetic problem generation statistics.}
\end{table}

\section{Process Reward Model}

\subsection{Motivation}

Outcome reward models (ORMs) assign a single score to the complete solution.
This signal is sparse for long multi-step problems and provides no information
about which steps are correct. Process reward models (PRMs) evaluate each reasoning
step, providing dense supervision:

\begin{equation}
r_\text{PRM}(s_1, s_2, \ldots, s_T) = \prod_{t=1}^T p(\text{step}_t \text{ correct} \mid s_1, \ldots, s_t)
\end{equation}

\subsection{PRM Training Data: Process Supervision}

We collect process supervision data via two methods:

\textbf{Human annotation}: 12K solutions to MATH and AMC problems annotated
step-by-step (correct/incorrect/uncertain) by mathematically trained annotators.
Inter-annotator agreement: 91.3\% on step-level labels.

\textbf{Monte Carlo sampling}: For a solution prefix $(s_1, \ldots, s_t)$, we estimate
step correctness by sampling $K=128$ completions and computing the fraction that
reach the correct final answer:

\begin{equation}
p(\text{step}_t \text{ correct}) \approx \frac{1}{K}\sum_{k=1}^K \mathbf{1}[\text{completion}_k \text{ correct}]
\end{equation}

\begin{algorithm}[H]
\caption{Monte Carlo PRM Data Generation}
\begin{algorithmic}[1]
\REQUIRE Problem set $\mathcal{P}$, policy $\pi$, completions per step $K=128$
\ENSURE Step-label dataset $\mathcal{D}_\text{prm}$
\FOR{each problem $p \in \mathcal{P}$}
  \STATE $S \leftarrow \pi.\text{solve}(p)$ \COMMENT{Generate full solution $S = (s_1, \ldots, s_T)$}
  \FOR{step $t = 1 \ldots T$}
    \STATE $\text{completions} \leftarrow [\pi(p, s_{1:t}) \text{ for } k=1\ldots K]$
    \STATE $\hat{q}_t \leftarrow \frac{1}{K}\sum_k \mathbf{1}[\text{correct}(\text{completion}_k)]$
    \STATE $\mathcal{D}_\text{prm} \leftarrow \mathcal{D}_\text{prm} \cup \{(p, s_{1:t}, \hat{q}_t)\}$
  \ENDFOR
\ENDFOR
\RETURN $\mathcal{D}_\text{prm}$
\end{algorithmic}
\end{algorithm}

\subsection{PRM Architecture and Training}

The PRM is a Zen-7B model with a classification head on each reasoning step token.
Training uses a binary cross-entropy loss with soft labels from MC sampling:

\begin{equation}
\mathcal{L}_\text{PRM} = -\sum_{t} \left[\hat{q}_t \log p_\phi(c_t) + (1-\hat{q}_t)\log(1-p_\phi(c_t))\right]
\end{equation}

We train on 800K step-labeled solutions (600K synthetic + 200K human-annotated).

\subsection{PRM-Guided Search}

At inference, we use beam search guided by the PRM:

\begin{equation}
\text{score}(s_{1:T}) = \min_t p_\phi(c_t \mid s_{1:t}) \quad (\text{pessimistic PRM})
\end{equation}

The pessimistic aggregation (minimum over steps) prioritizes solutions with no
weak steps over solutions with some strong steps and some weak ones.

\section{Hybrid Verification}

\subsection{Symbolic Verification (SymPy)}

For algebraic expressions, equations, and numerical answers, we use SymPy to
verify equivalence:

\begin{algorithm}[H]
\caption{Symbolic Verification}
\begin{algorithmic}[1]
\REQUIRE Predicted answer $\hat{a}$, ground truth $a^*$
\ENSURE Boolean: correct or not
\STATE Parse $\hat{a}$ and $a^*$ into SymPy expressions
\IF{simplify$(\hat{a} - a^*) = 0$}
  \RETURN True
\ELSIF{simplify$(\hat{a} - a^*) \neq 0$}
  \RETURN False
\ELSE
  \STATE \RETURN NeuralVerify$(\hat{a}, a^*)$
\ENDIF
\end{algorithmic}
\end{algorithm}

Symbolic verification handles polynomial equations, trigonometric identities, and
rational expressions. It fails on open-form answers (proofs, geometric arguments,
combinatorial constructions), which are routed to neural verification.

\subsection{Neural Verification}

For non-algebraic answers, a Zen-32B model trained on 50K verified proof/answer
pairs assesses correctness. The verifier is given both the predicted answer and
the ground truth and produces a calibrated correctness probability.

\section{Experiments}

\subsection{Benchmark Results}

\begin{table}[H]
\centering
\begin{tabular}{lccccc}
\toprule
\textbf{Model} & \textbf{MATH} & \textbf{GSM8K} & \textbf{AIME 2024} & \textbf{AMC 2023} & \textbf{FrontierMath} \\
\midrule
Zen-7B & 71.4 & 95.8 & 38.3 & 74.2 & 28.1 \\
Zen-7B + PRM & 76.2 & 96.8 & 40.1 & 77.4 & 29.8 \\
Zen-32B & 78.8 & 97.2 & 44.3 & 81.2 & 33.4 \\
Zen-32B + PRM & \textbf{83.4} & \textbf{97.9} & \textbf{48.2} & \textbf{84.7} & \textbf{36.9} \\
\bottomrule
\end{tabular}
\caption{Mathematical reasoning benchmark results. PRM adds 4--8 points across benchmarks.}
\end{table}

Note: The headline Zen-7B result of 73.2\% MATH cited in the abstract reflects the
production model with best-of-32 sampling; single-sample accuracy is 71.4\%.

\subsection{MATH Category Breakdown}

\begin{table}[H]
\centering
\begin{tabular}{lcc}
\toprule
\textbf{Category} & \textbf{Zen-7B (\%)} & \textbf{Zen-7B + PRM (\%)} \\
\midrule
Algebra & 84.3 & 88.1 \\
Counting \& Probability & 71.2 & 76.4 \\
Geometry & 62.4 & 67.8 \\
Number Theory & 78.3 & 82.1 \\
Precalculus & 66.1 & 71.2 \\
Intermediate Algebra & 69.8 & 74.3 \\
\midrule
Overall & 71.4 & 76.2 \\
\bottomrule
\end{tabular}
\caption{MATH benchmark by category. Geometry benefits most from PRM guidance.}
\end{table}

\subsection{Effect of Sampling at Test Time}

\begin{table}[H]
\centering
\begin{tabular}{lccccc}
\toprule
\textbf{Samples} & \textbf{1} & \textbf{4} & \textbf{16} & \textbf{32} & \textbf{64} \\
\midrule
No PRM (majority vote) & 71.4 & 74.1 & 76.3 & 77.2 & 77.8 \\
PRM (best-of-$k$) & 71.4 & 74.8 & 77.3 & 78.1 & 79.0 \\
\bottomrule
\end{tabular}
\caption{MATH accuracy vs. number of samples (Zen-7B). PRM reranking outperforms majority vote.}
\end{table}

\subsection{PRM Calibration}

The PRM step correctness scores are well-calibrated: across 50K held-out steps,
steps scored $<0.1$ have 6.2\% actual correctness rate; steps scored $>0.9$ have
91.4\% actual correctness rate. Calibration error (ECE) = 0.043.

\section{Analysis}

\subsection{Error Taxonomy}

We manually analyzed 200 errors by Zen-7B on MATH Level 5 problems:

\begin{table}[H]
\centering
\begin{tabular}{lrc}
\toprule
\textbf{Error Type} & \textbf{Count} & \textbf{PRM Detectable?} \\
\midrule
Arithmetic computation error & 61 & 78\% \\
Wrong approach / setup & 44 & 52\% \\
Algebraic manipulation error & 38 & 71\% \\
Missing case / edge case & 29 & 31\% \\
Incorrect interpretation & 18 & 18\% \\
Copy/transcription error & 10 & 90\% \\
\midrule
Total & 200 & \\
\bottomrule
\end{tabular}
\caption{Error taxonomy for Zen-7B on MATH Level 5. PRM detects most computational errors.}
\end{table>

The PRM is most effective at detecting arithmetic and algebraic errors (step produces
numerically inconsistent result) and least effective at detecting missing cases or
misinterpretation (where each individual step appears locally valid).

\subsection{Olympiad-Level Performance}

On AIME 2024 (30 problems), Zen-32B with PRM-guided search and $n=128$ samples
achieves 48.2\% (14.5/30 problems). Human AIME competitors typically score 8--15
in competitive settings.

On FrontierMath (research-level problems from active mathematics research), 36.9\%
represents a significant capability frontier. Many FrontierMath problems require
theorem-proving capabilities beyond current neural approaches.

\section{Conclusion}

Mathematical reasoning in Zen models benefits substantially from three components:
domain-rich pretraining data (120B tokens), process reward models that provide
step-level guidance (+4--8 MATH points), and hybrid verification enabling reliable
self-evaluation. The combination yields 83.4\% MATH, 97.9\% GSM8K, and 48.2\% AIME
for Zen-32B. Open problems include proof-level verification for advanced mathematics
and step-level data collection at Olympiad difficulty.

\bibliographystyle{plain}
\begin{thebibliography}{99}
\bibitem{lightman2023prm} Lightman et al. (2023). Let's Verify Step by Step. \textit{arXiv:2305.20050}.
\bibitem{hendrycks2021math} Hendrycks et al. (2021). Measuring Mathematical Problem Solving With the MATH Dataset. \textit{NeurIPS}.
\bibitem{cobbe2021gsm8k} Cobbe et al. (2021). Training Verifiers to Solve Math Word Problems. \textit{arXiv:2110.14168}.
\bibitem{wang2023prm800k} Wang et al. (2023). Math-Shepherd: Verify and Reinforce LLMs Step-by-Step without Human Annotations. \textit{arXiv:2312.08935}.
\bibitem{trinh2024alphageometry} Trinh et al. (2024). Solving olympiad geometry without human demonstrations. \textit{Nature}.
\end{thebibliography}

\end{document}
